%% Kelp Segmentation Datasets Survey
%% Comprehensive Report for PhD Research
%% Last Updated: March 2026

\documentclass[11pt]{article}
\usepackage[utf8]{inputenc}
\usepackage[T1]{fontenc}
\usepackage[margin=1in]{geometry}
\usepackage{booktabs}
\usepackage{xcolor}
\usepackage{enumitem}
\usepackage{hyperref}
\usepackage{titlesec}
\usepackage{longtable}
\usepackage{array}
\usepackage{makecell}

% Colors
\definecolor{available}{RGB}{39,174,96}
\definecolor{contact}{RGB}{243,156,18}
\definecolor{unavailable}{RGB}{231,76,60}

% Section formatting
\titleformat{\section}{\large\bfseries\color{blue!70!black}}{}{0em}{}[\titlerule]
\titleformat{\subsection}{\normalsize\bfseries}{}{0em}{}

\title{\textbf{Kelp Segmentation Datasets}\\[0.5em]
\large Data Survey for LA-Based Kelp Detection Research}
\author{Ekaba Bisong}
\date{March 2026}

\begin{document}

\maketitle
\tableofcontents


%% ============================================================
\section{Executive Summary}
%% ============================================================

This document provides a survey of kelp segmentation datasets for training and validating Learning Automata-based kelp detection systems. The focus is on:
\begin{itemize}
  \item Giant kelp (\textit{Macrocystis pyrifera}) -- California, Baja California
  \item Bull kelp (\textit{Nereocystis luetkeana}) -- Pacific Northwest
  \item Satellite imagery from Landsat and Sentinel-2
\end{itemize}

\subsection*{Dataset Availability Summary}

\begin{center}
\begin{tabular}{llcl}
\toprule
\textbf{Dataset} & \textbf{Type} & \textbf{Status} & \textbf{Size} \\
\midrule
Figshare Kelp (Mask R-CNN) & Labeled imagery & \textcolor{available}{Downloaded} & 73 MB \\
Floating Forests & Citizen science labels & \textcolor{contact}{Contact Required} & 5,635 images \\
Hakai Institute & Labeled aerial/satellite & \textcolor{contact}{Contact Required} & Unknown \\
\bottomrule
\end{tabular}
\end{center}

%% ============================================================
\section{Downloaded Datasets}
%% ============================================================

\subsection{1. Figshare Mask R-CNN Kelp Dataset}

\begin{tabular}{p{4cm}p{10cm}}
\toprule
\textbf{Attribute} & \textbf{Details} \\
\midrule
Location & \texttt{data/figshare\_kelp/} \\
Size & 73 MB (432 images) \\
Source & Figshare (DOI: 10.6084/m9.figshare.19935869) \\
Region & Southern California and Baja California \\
Species & Giant kelp (\textit{Macrocystis pyrifera}) \\
Satellite & Landsat Thematic Mapper \\
Resolution & 30 m \\
Format & JPEG images + VIA JSON annotations \\
Labels & Polygon masks for kelp canopy \\
Splits & Train: 323, Val: 30, Test: 79 \\
Performance & Dice: 0.93, Jaccard: 0.87 \\
Code & \url{https://github.com/jorgeassis/maskRCNN} \\
Citation & Marquez et al. (2022) Scientific Reports \\
\bottomrule
\end{tabular}


%% ============================================================
\section{Datasets Requiring Contact}
%% ============================================================

The following high-quality dataset may be available upon request:

\subsection{1. Floating Forests (DrivenData Source)}

\begin{tabular}{p{4cm}p{10cm}}
\toprule
\textbf{Attribute} & \textbf{Details} \\
\midrule
Project & Floating Forests (Zooniverse citizen science) \\
Labels & $>$1 million citizen classifications \\
Images & 5,635 labeled Landsat images (DrivenData competition) \\
Regions & Falkland Islands (100\%), California, Tasmania \\
Validation & Each image shown to 15 users for consensus \\
Status & DrivenData competition ended (Feb 2024); data removed from site \\
\bottomrule
\end{tabular}

\vspace{0.5em}
\textbf{Contact Information:}

\textbf{Dr.\ Kyle Cavanaugh} (UCLA Geography)
\begin{itemize}[nosep]
  \item Email: \texttt{kcavanaugh@geog.ucla.edu}
  \item Web: \url{https://geog.ucla.edu/person/kyle-cavanaugh/}
\end{itemize}

\textbf{Dr.\ Jarrett Byrnes} (UMass Boston Biology)
\begin{itemize}[nosep]
  \item Email: \texttt{Jarrett.Byrnes@umb.edu}
  \item Lab: \url{https://byrneslab.net/}
\end{itemize}

\textbf{Why valuable:} Original labeled dataset with 5,635 training images and validated citizen science labels. Includes cross-hemisphere data (Falkland Islands).

\subsection{2. Hakai Institute (BC, Canada)}

\begin{tabular}{p{4cm}p{10cm}}
\toprule
\textbf{Attribute} & \textbf{Details} \\
\midrule
Institution & Hakai Institute, British Columbia \\
Region & BC Central Coast (Canada) \\
Species & Giant kelp (\textit{Macrocystis}), Bull kelp (\textit{Nereocystis}) \\
Data Types & Field data, aerial imagery, satellite-derived products \\
ML Training & Semantic segmentation labels (NPZ format) \\
GEEK Tool & Landsat kelp extent 1984--2019 (Google Earth Engine) \\
GitHub & \url{https://github.com/HakaiInstitute/hakai-ml-train} \\
\bottomrule
\end{tabular}

\vspace{0.5em}
\textbf{Available Resources:}
\begin{itemize}[nosep]
  \item Kelp field data matched to satellite imagery for spectral signatures
  \item Annual kelp extent rasters (1984--2019) from Landsat
  \item High-resolution airborne coastal observatory data
\end{itemize}

\textbf{Contact:} Data catalogue at \url{https://catalogue.hakai.org/}

\textbf{Why valuable:} Species-specific labels (Macrocystis vs Nereocystis), field-validated ground truth, and existing ML training infrastructure.

%% ============================================================
\section{Datasets NOT Useful for Kelp Research}
%% ============================================================

The following datasets were evaluated but are \textbf{not relevant} for kelp forest research:

\begin{center}
\begin{tabular}{lll}
\toprule
\textbf{Dataset} & \textbf{Reason Not Useful} & \textbf{Actual Target} \\
\midrule
USF Floating Algae & Different organisms & Sargassum, Cyanobacteria, Ulva \\
Marine Debris Data & Wrong task & Plastic, driftwood, debris detection \\
IMO Macroalgae & Different species/habitat & European intertidal algae (not kelp) \\
\bottomrule
\end{tabular}
\end{center}

\textbf{Key distinction:} Kelp (Macrocystis, Nereocystis) are subtidal canopy-forming species attached to the seafloor. The above datasets target floating algae, debris, or intertidal species.

\end{document}
